\section{财产犯罪 - 抢劫罪}
\begin{enumerate}
  \item 保护法益。人身 + 财产。
  \item \zd{非法占有目的 + 暴力、胁迫手段 \yjt 压制反抗 \yjt 因无法反抗而放弃财物 \yjt 因此取得财物}。
        % https://q.uiver.app/?q=WzAsNCxbMCwwLCLpnZ7ms5XljaDmnInnm67nmoQgKyDmmrTlipsiXSxbMSwwLCLljovliLblj43mipciXSxbMiwwLCLlm6Dml6Dms5Xlj43mipdcXFxc6ICM5pS+5byD6LSi54mpIl0sWzMsMCwi5Y+W5b6X6LSi54mpIl0sWzAsMV0sWzEsMl0sWzIsMywi5pei6YGC5p2h5Lu2XFxcXOWboOaenOaApyIsMV0sWzAsMywi5pe256m657Sn5a+G5oCnIiwxLHsiY3VydmUiOjJ9XV0=
\[
  \begin{tikzcd}
    \txt{非法占有目的 + 暴力} & \txt{压制反抗} & \txt{因无法反抗\\而放弃财物} &&& \txt{取得财物}
    \arrow[from=1-1, to=1-2]
    \arrow[from=1-2, to=1-3]
    \arrow["\txt{既遂条件\\因果性}"{description}, from=1-3, to=1-6]
    \arrow["\txt{时空紧密性}"{description}, curve={height=36pt}, from=1-1, to=1-6]
  \end{tikzcd}
\]

        \begin{enumerate}
          \item 成立条件。\zd{行为与目的同时存在原则}。
                \begin{enumerate}
                  \item 行为对象。主人、占有者、占有辅助者、（临时）\zd{保护者}。
                        \begin{enumerate}
                          \item \al{甲乙盗窃网吧熟睡丙，老板上前阻止被威胁}。
                          \item \al{甲盗窃熟睡中的乙，小孩丙嘲笑之，甲威胁丙}。
                          \item \al{甲看到商店门口停着电动车，乙站在旁边，以为乙是车主，暴打一顿，骑走电车 \cuo}。
                        \end{enumerate}
                  \item 暴力、胁迫。\zd{以恶害相通告，产生恐惧心理，足以压制对方反抗}。2、3 步因果关系。
                        \begin{enumerate}
                          \item \al{甲乙抢劫银行，进门后说不给钱我就引爆，实际是牛油果}。
                          \item \al{小偷跑，主人追，揉了一团报纸砸向主人}。
                          \item \al{甲抢劫乙女，乙女追赶，甲脱下裤子}。
                          \item 昏醉抢劫、拘禁型。
                                \al{甲乙共谋诱骗丙参赌，下药趁其不清醒合伙赢钱，结果剂量太大，刚上牌桌丙就睡去，甲乙趁机将丙的钱包掏空离去 \yjt 抢劫}。
                                \al{甲乙丙一起赌博，丙熬夜太晚睡过去了，甲乙趁机将丙的钱包掏空离去 \yjt 盗窃}。非法占有目的产生的阶段不同。
                        \end{enumerate}
                  \item 暴力行为、取财行为要有\zd{时空紧密性}。当场性。1、4 步。
                        \begin{enumerate}
                          \item \al{高铁得知是富婆，杀人取钱}。
                          \item \al{得知是富婆，知道家里有钱，杀人去杭州取财}。
                        \end{enumerate}
                \end{enumerate}
          \item 既遂。无法反抗被迫放弃 \yjt \zd{取得财物}（第二档，数额不大）要有因果关系。3、4 步。
                \begin{enumerate}
                  \item \al{甲抢劫乙，乙跑途中钱包掉落}。乙并非无法反抗放弃财物。
                  \item 既遂 2。\zd{造成轻伤以上}（司法解释）。
                \end{enumerate}
        \end{enumerate}
  \item \zd{事后转化抢劫}。拟制。
        \begin{enumerate}
          \item \zd{14 - 16 岁的人不定转化抢劫 \yjt 致人重伤、死亡 \yjt 故意伤害、故意杀人}。
          \item 3 个轻罪。盗窃、诈骗、抢夺。\zd{行为具有财产犯罪的属性 + 符合这 3 个罪的构成要件}。
          \item 3 个目的。\zd{窝藏赃物、抗拒抓捕、毁灭罪证}。\al{受害人随后追赶}。
                \begin{enumerate}
                  \item \al{公交车抢包，到站跳下车，被害人破口大骂，跳上车扇两耳光 \cuo}。
                  \item \al{甲抢劫乙项链到手，以为是假的丢给乙，扇了两耳光，骂道出门也不带好的 \cuo}。
                \end{enumerate}
          \item 实行行为。\zd{当场使用暴力或以暴力相威胁，足以压制对方反抗（程度）}。
                \begin{enumerate}
                  \item 针对的主体是人。被害人、抓捕人、妨碍人。\al{无辜路人甲 \cuo}。
                  \item 当场性。时间概念。前行为和后面的暴力行为要有\zd{时间紧密性}。喜欢考时空错位，\zd{把现场当当场}。
                        \begin{enumerate}
                          \item 空间上是现场，时间上不是当场。\al{团伙盗窃别墅完事后做饭吃，次日主人回来，把主人打了}。
                          \item 空间上不是同一现场，时间上是当场。\al{持续追捕}。
                        \end{enumerate}
                  \item 以暴力相威胁。针对的主体仍然是人。\al{如果不停止抓捕，就会对其实施暴力}。
                  \item 为抗拒抓捕而使用暴力\zd{要不要主客观一致}。\al{甲抢夺后逃跑，乙追赶，甲抗拒追捕开枪不慎打死丙}。
                        \begin{enumerate}
                          \item \zd{多数说}。要求 \yjt 不构成转化。
                          \item 少数说。只要求主观，客观不要求 \yjt 构成。
                          \item \sjt 不影响最后结论 \yjt 转化抢劫 \yjt 抢劫罪致人重伤/死亡。
                        \end{enumerate}
                  \item 既遂。造成他人轻伤。
                \end{enumerate}
          \item vs 临时起意、直接升级。\zd{核心。行为与目的同时存在原则}。
                \begin{enumerate}
                  \item 使用暴力的目的不同。转化抢劫目的是上面 3 个（因为已取得财物 \yjt 不要求既遂的平稳状态）。后者是为了继续取得财物。
                  \item \al{伤害、强奸后后临时起意，没有知觉/不能反抗、不敢反抗 \yjt 拿走/暴力、胁迫手段 \yjt 盗窃/抢劫。数罪并罚}。
                  \item \al{甲为强奸乙女，将其打倒重伤在地，但仍清醒。甲看到掉落在身旁的钱包，捡起离去 \yjt 不知反抗后 \yjt 未实施抢劫行为 \yjt 盗窃}。
                  \item \al{甲飞车抢夺乙的背包，乙使劲抓住背包带，甲于是加速行驶，乙被拖行 10m 多后松手。直接升级为抢劫，不属于转换型}！
                \end{enumerate}
          \item 共犯问题。
                \begin{enumerate}
                  \item 共犯与身份。犯盗窃者、……才能构成\zd{转化抢劫 \yjt 是一种真正身份犯(含实行者、帮助者)}。
                        \begin{enumerate}
                          \item \al{没有身份，不能成为转化抢劫的实行犯 \yjt
                                  乙没有 \yjt
                                  教唆犯要求存在实行犯 \yjt
                                  甲不构成转化抢劫的教唆犯 \yjt
                                  拟制甲为转化抢劫的间接正犯 \yjt
                                  乙构成转化抢劫的帮助犯 \yjt
                                  甲乙构成转化抢劫的共同犯罪}。
                          \item \al{甲单独盗窃既遂逃离现场，丙追捕，乙中途参与，教唆甲打丙}。
                          \item \al{甲\xhx{让乙帮忙}断后，乙绊倒丙 \yjt 甲转化抢劫间接正犯、故意伤害教唆犯。乙转化抢劫帮助犯、故意伤害罪正犯 \yjt 想象竞合择一重}。
                          \item \al{甲\xhx{骗乙}有人追杀他，乙不知情绊倒丙 \yjt 甲转化抢劫间接正犯。乙假想防卫}。
                          \item \al{乙暗中相助绊倒丙 \yjt 乙窝藏、故意伤害。甲盗窃既遂}。
                        \end{enumerate}
                  \item \zd{实行过限，谁暴力谁负责}。只考过第三种。\al{甲乙共同盗窃，乙望风}。
                        \begin{enumerate}
                          \item \al{甲绊倒丙 \yjt 甲转化。If(乙不知道) \yjt 盗窃，帮助犯}。
                          \item \al{丙抓乙，乙绊倒丙}。
                          \item \al{丙追甲，乙绊倒丙帮甲断后，\xhx{甲不知} \yjt 乙具有身份（帮助者）转化抢劫。甲盗窃}。
                        \end{enumerate}
                \end{enumerate}
        \end{enumerate}
  \item \zd{法定刑升格条件}。
        \begin{enumerate}
          \item 入户抢劫。入户盗窃转为入户抢劫。
                \begin{enumerate}
                  \item 入前就要求有非法目的。\al{甲送乙女回家，乙要求甲上去喝咖啡，甲抢劫之}。引狼入室。
                  \item 户要求生活场所（看当时状态）。\al{租房卖淫，嫖客完事后抢劫}。
                  \item \zd{暴力行为要求一定在户内}。
                        \begin{enumerate}
                          \item \al{甲在户外开枪}。
                          \item \al{甲盗窃乙家被发现，乙追至户外，甲痛打一顿}。
                        \end{enumerate}
                  \item 和转化抢劫结合一起考。\al{甲盗窃乙家过程中被发现，抗拒抓捕对乙暴力}。
                \end{enumerate}
          \item 在公共交通工具上抢劫。
                \begin{enumerate}
                  \item 公共交通工具缩小解释。不包括出租车摩的，得大中型。
                  \item 上。\zd{不要求上去，这里是针对的意思}。
                  \item 暴力行为在车上。因为在车上才可能对不特定多数人产生危险，在车外那就是针对特定人了。
                \end{enumerate}
          \item 抢劫银行或其他金融机构（\zd{经营资金}）、多次(3 次以上)抢劫。
          \item \zd{抢劫 + 致人重伤/死亡 = 抢劫加重处罚}。\al{所以看到抢劫、故意伤害出现在一起一律不选}。
                \begin{enumerate}
                  \item 人的范围。被害人、抓捕人、妨碍人、\zd{打击错误第三人（含同伙）}。
                        \al{甲乙抢劫丙，丙呼喊反抗，甲拔刀刺向丙，丙闪躲，刺死了乙 \dui}。
                  \item 含转化抢劫的情形。
                  \item 区别人身类犯罪。\al{甲欲强奸乙女，丙阻止，甲重伤之。分别定，而不是强奸致人重伤、死亡}。
                  \item \zd{注意因果关系，介入因素}。
                \end{enumerate}
          \item 冒充军警人员抢劫。\yi 能使一般人相信。\er \zd{得是假警察}。
          \item 持枪抢劫。\zd{真枪但不要求实弹}。需显示 \hz 使用。
          \item 抢劫军用物资 \hz 抢险、救灾、救济物资。
        \end{enumerate}
  \item vs 绑架。区别在于行为时带的目的，两个罪可以竞合。\al{甲一把抢过乙的孩子，不给钱捅死他}。
\end{enumerate}
